\section{Related Work}

In the domain of PDE solvers, quantum algorithms can be broadly categorized into two primary approaches: (a) exclusively quantum algorithms and (b) hybrid quantum-classical methods. 

\subsection{Exclusively quantum algorithms}

Exclusively quantum algorithms for solving PDEs typically leverage unitary operations and quantum state encoding techniques to achieve computational advantages. In contrast to neural networks, exclusively quantum algorithms use fixed quantum circuit structures with no trainable parameters.
The celebrated Harrow-Hassidim-Lloyd (HHL) algorithm is an example of this category, which provides exponential speedups for solving linear systems under specific conditions~\cite{harrow2009quantum}. Extensions of the HHL algorithm have been applied to solve sparse linear systems derived from discretized linear~\cite{ambainis2012variable, somma2016quantum} or nonlinear~\cite{leyton2008quantum} PDEs, leveraging numerical techniques such as linear multistep method~\cite{berry2014high}, Chebyshev pseudospectral method~\cite{childs2020quantum}, and Taylor series approximations~\cite{berry2017quantum}. In fluid dynamics applications, lattice-gas quantum models have been examined for directly solving PDEs~\cite{yepez2002quantum,montanaro2016quantum} or PDEs derived from ODEs through methods such as the Quantum Amplitude Estimation Algorithm (QAEA)~\cite{oz2022solving,oz2023efficient}.

% Problem / Gap
Despite these advancements, purely quantum algorithms for solving PDEs encounter practical challenges such as requirements for quantum error correction~\cite{preskill1998reliable, terhal2015quantum, childs2021theory}, 
limitations in circuit depth~\cite{bravyi2005universal, montanaro2016quantum, boixo2018characterizing}, and bottlenecks in data encoding~\cite{aaronson2015read, ciliberto2018quantum, lloyd2013quantum}. These obstacles, among others, limit the near-term viability of exclusively quantum PDE solvers on today's noisy intermediate-scale quantum (NISQ) devices.

\subsection{Hybrid quantum-classical approaches} 

Hybrid quantum-classical approaches have emerged as promising solutions to overcome the limitations of exclusively quantum algorithms. These methods combine the robustness, maturity, and stability of classical computation with the enhanced capacity, expressivity, and inherent parallelism of quantum systems. While classical components can efficiently handle pre-/post-processing, optimization, and error stabilization, quantum layers can capture intricate patterns, solve sub-problems, and accelerate specialized computations.

Within this hybrid paradigm, Quantum Neural Networks (QNNs) have demonstrated universal approximation capabilities for continuous functions when implemented with sufficient depth, width, and well-designed quantum feature maps~\cite{schuld2019quantum, lloyd2020quantum, beer2020training}. For instance, Salinas et al.\cite{PhysRevA.104.012405} showed that single-qubit networks with iterative input reuploading can approximate bounded complex functions, while Goto et al.\cite{goto2021universal} established universal approximation theorems for quantum feedforward networks. Moreover, Schuld et al.\cite{schuld2021effect} revealed that the outputs of parameterized quantum circuits (with suitable data encoding and circuit depth) can be expressed as truncated Fourier series, thereby providing theoretical support for the universal approximation ability of QNNs. Furthermore,  Liu et al.\cite{liu2021rigorous} provided formal evidence of quantum advantage in function approximation. These developments not only support applications in solving PDEs and optimization tasks but also extend to other computational tasks such as quantum chemistry~\cite{cao2019quantum}, materials science~\cite{bauer2020quantum, nachman2021quantum}, and pattern recognition~\cite{trugenberger2002quantum} to name a few.

A well-known example of this integrated approach is the Variational Quantum Algorithms (VQA), also known as Variational Quantum Circuits (VQC). These algorithms employ parameterized quantum circuits with tunable gates optimized through classical feedback loops. VQA implementations for solving PDEs include the Variational Quantum Eigensolver (VQE)\cite{peruzzo2014variational,mcclean2016theory}, Variational Quantum Linear Solver (VQLS)\cite{bravo2023variational}, and Quantum Approximate Optimization Algorithm (QAOA)\cite{albino2022solving}, as well as emerging methods like Differentiable Quantum Circuits (DQCs)\cite{kyriienko2021solving,bultrini2023mixed}, quantum kernel methods~\cite{paine2023quantum}, and Quantum Nonlinear Processing Algorithms (QNPA)~\cite{lubasch2020variational}. 


The aforementioned methods primarily utilize qubit-based quantum computing, which represents the predominant paradigm in quantum computing literature. However, continuous-variable quantum computing (CVQC)~\cite{braunstein2005quantum} offers an alternative computational model that leverages continuous degrees of freedom, such as quadratures and electromagnetic field momentum. For instance, in~\cite{markidis2022physics}, Gaussian and non-Gaussian gates were used to manage continuous degrees of freedom, offering enhanced scalability and expressiveness for continuous systems. For consistency, we refer to qubit-based models as discrete-variable (DV) models throughout this paper.

Recently, VQA methods have been extended to Physics-informed quantum neural networks (PIQNN) \cite{setty2023physics, dehaghani2024hybrid, trahan2024quantum,sedykh2024hybrid, leong2025hybrid}, broadening their applicability across physics and engineering domains. These hybrid methods offer several distinct advantages when implemented within the PINN framework. 

Firstly, these models enable the flexible configuration of inputs and outputs independent of qubits/qumodes counts or quantum layer depth, making them particularly effective for scenarios where the input dimension (e.g., temporal and spatial coordinates) differs from output dimensions (e.g., dependent variables like velocity, pressure, and force).

Secondly, hybrid approaches facilitate increased trainable parameter counts without necessitating deeper quantum circuits, additional qubits, or higher cutoff dimensions in continuous-variable (CV) quantum systems. These factors are particularly crucial in the NISQ era, where hardware constraints impose strict limits on circuit depth and coherence times, making deeper quantum networks computationally expensive and impractical~\cite{preskill2018quantum, bharti2022noisy}. 

Thirdly, these qualities are especially advantageous for PINNs, where shallow network architectures are favored to optimize the balance between computational efficiency and solution accuracy~\cite{krishnapriyan2021characterizing, wang2022approximation, shin2023error, farea2024:arXiv}

Finally, a substantial benefit of such hybrid models is the integration of automatic differentiation within both classical and quantum networks. This capability, supported by contemporary software packages like PennyLane~\cite{bergholm2018pennylane} and TensorFlow Quantum~\cite{broughton2020tensorflow}, enables the automatic computation of derivatives as variables are directly incorporated into the computational graph. This approach simplifies PINN implementation and facilitates the efficient handling of complex physics-informed loss functions and higher-order derivatives.

The work presented here is related to the approaches in~\cite{sedykh2024hybrid,dehaghani2024hybrid,trahan2024quantum,leong2025hybrid}, but with several key differences. While the HQPINN~\cite{sedykh2024hybrid} focuses specifically on computational fluid dynamics in 3D Y-shaped mixers using a fixed quantum depth with an infused layer structure, the QPINN~\cite{dehaghani2024hybrid} integrates a dynamic quantum circuit combining Gaussian and non-Gaussian gates with classical computing methodologies within a PINN framework to solve optimal control problems. Similarly, the model proposed in~\cite{trahan2024quantum} investigates both purely quantum and hybrid PINNs, highlighting parameter efficiency but relying on limited PQC networks and addressing primarily simple PDEs. Recently proposed HQPINN~\cite{leong2025hybrid}, inspired by Fourier neural operator (FNO)~\cite{li2020fourier}, introduces a parallel hybrid architecture to separately model harmonic and non-harmonic features before combining them for output.